\chapter{Criterio de clausura editorial provisional del sistema}

\section{Clausura editorial provisional}
\begin{definition}
Se llama clausura editorial provisional del sistema al estado disciplinado en el que MADELEINE puede declararse suficientemente organizado, estabilizado y publicable en su forma actual, sin que ello implique agotamiento absoluto de toda expansión futura del corpus.
\end{definition}

\begin{observation}
La clausura editorial provisional no equivale a cierre metafísico ni a final definitivo del sistema. Su función es fijar una forma de publicación y presentación suficientemente madura para sostener una salida externa sobria.
\end{observation}

\section{Condición mínima de clausura editorial provisional}
\begin{definition}
Existe clausura editorial provisional cuando pueden declararse simultáneamente:
\begin{enumerate}
  \item una arquitectura legible entre Tomo I, Tomo II y Tomo III;
  \item una distinción suficiente entre núcleo, canon intermedio, metodología, capa pública y material reservado;
  \item una política editorial sobria de publicación;
  \item un cierre operativo suficiente de Tomo III;
  \item una forma externa del sistema que no contradiga la jerarquía metodológica interna.
\end{enumerate}
\end{definition}

\begin{observation}
No se exige todavía cancelación de futuras adendas, extensiones o revisiones. Basta con que el sistema ya pueda sostener una forma editorial coherente en su fase presente.
\end{observation}

\section{Clausura y provisionalidad}
\begin{definition}
La provisionalidad de la clausura editorial significa que:
\begin{enumerate}
  \item el sistema puede publicarse sin quedar clausurado ontológicamente;
  \item futuras ampliaciones deben entrar como continuidad disciplinada y no como destrucción del cierre ya alcanzado;
  \item la arquitectura publicada conserva prioridad sobre expansiones posteriores mientras no se reabra expresamente el marco editorial.
\end{enumerate}
\end{definition}

\begin{observation}
Esta cláusula protege al corpus frente a dos extremos: la ilusión de final absoluto y la disolución permanente de todo cierre posible.
\end{observation}

\section{Clausura y jerarquía entre tomos}
\begin{definition}
Toda clausura editorial provisional debe preservar:
\begin{enumerate}
  \item Tomo I como fijación del núcleo y del canon basal;
  \item Tomo II como canon intermedio y organización operativa;
  \item Tomo III como capa metodológica, curatorial y de proyección disciplinada.
\end{enumerate}
\end{definition}

\begin{observation}
La clausura editorial no debe colapsar las funciones de los tomos ni convertir el sistema en una sola masa textual indiferenciada.
\end{observation}

\section{Condición mínima de no clausura}
\begin{definition}
No debe declararse clausura editorial provisional cuando la forma actual del sistema siga dependiendo de vacíos estructurales, ambigüedades fuertes entre tomos o reglas editoriales insuficientes para sostener una publicación coherente.
\end{definition}

\begin{observation}
La no clausura editorial provisional no invalida el sistema. Puede significar solo que todavía conviene consolidar arquitectura, política editorial o mediaciones externas antes de fijar la fase presente.
\end{observation}

\section{Cláusula de alcance}
\begin{observation}
El criterio de clausura editorial provisional del sistema fija la primera condición de cierre publicable de MADELEINE como corpus organizado. No constituye todavía una clausura absoluta de toda expansión futura ni una prohibición de reabrir el sistema bajo un nuevo ciclo disciplinado.
\end{observation}

